\documentclass[journal]{IEEEtran}

\usepackage{cite}
\usepackage{amsmath,amssymb}
\usepackage{graphicx}
\usepackage{textcomp}
\usepackage{xcolor}
\usepackage{booktabs}
\usepackage{multirow}
\usepackage{array}
\usepackage[hidelinks]{hyperref}

\begin{document}

\title{CIDAS: A Context-Informed Dependency Analysis System for Real-Time npm Supply-Chain Defense}

\author{Gayathri~Sunil~Nambiar,~Nitya~Sharma,~Sanvi~H~S,~Yuva~T,~and~Sriram~Velumuri
\thanks{G. S. Nambiar is with the Department of Computer Science (Data Science), RV College of Engineering, Bengaluru 560059, Karnataka, India (e-mail: gayathri.sunil@rvce.edu.in).}
\thanks{N. Sharma and S. Velumuri are with the Department of Artificial Intelligence and Machine Learning, RV College of Engineering, Bengaluru 560059, Karnataka, India (e-mail: \{nitya.sharma, sriram.velumuri\}@rvce.edu.in).}
\thanks{S. H. S and Y. T are with the Department of Computer Science and Engineering, RV College of Engineering, Bengaluru 560059, Karnataka, India (e-mail: \{sanvi.hs, yuva.t\}@rvce.edu.in).}
\thanks{Corresponding author: Gayathri Sunil Nambiar (e-mail: gayathri.sunil@rvce.edu.in).}}

\markboth{Nambiar \MakeLowercase{\textit{et al.}}: CIDAS: Context-Informed Dependency Analysis System}{}

\maketitle

\begin{abstract}
The npm registry---the world's largest open-source package repository---is an attractive target
for software supply-chain attacks. Malicious actors exploit developer trust through typosquatting, malicious
post-install scripts, maintainer account takeovers, and an emerging class of threat in which large language
model coding assistants hallucinate package names that adversaries subsequently register as weaponized
packages, a phenomenon termed slopsquatting. Existing defenses are predominantly reactive, scanning for
threats after packages have already been installed and their scripts executed. We present CIDAS (Context-
Informed Dependency Analysis System), a real-time, pre-install screening framework that intercepts
every npm install invocation through a transparent cross-platform shim and evaluates each candidate
dependency through four coordinated analysis pillars: Contextify (semantic project-context fit via sentence-
transformer embeddings), Sentinel (registry reputation analysis, curated supply-chain incident matching,
edit-distance and affix-canonicalized typosquat detection corroborated against reputation disparity, and
live OSV.dev advisory lookup for artificial-intelligence-suggested packages), Shield (static and abstract
syntax tree-based lifecycle-script and tarball scanning with manifest-gated cross-version capability diffing),
and a two-stage Aggregator that produces an interpretable 0--100 risk score mapped to ALLOW, WARN, or
BLOCK verdicts. Supporting subsystems include schema-validated policy-as-code governance, an HMAC-protected
per-machine trust list, an append-only audit log, a two-tier SQLite cache for daemon-outage resilience, and a
Kullback--Leibler divergence-based concept-drift monitor operating on genuine per-pillar score history. Evaluated
against labelled corpora of benign, typosquatted, hallucinated, and malicious packages (179 records,
including 17 real 2025--2026 supply-chain incidents added to the malicious corpus after the original
evaluation) with 548 automated tests (449 daemon, 99 extension), CIDAS achieves a pooled (aggregate) F1 of
0.93 and a category-balanced macro-F1 of 0.92, with sub-50\,ms cache-hit latency and single-digit-second
cache-miss latency on commodity hardware against live registry conditions,
without requiring cloud infrastructure or changes to existing developer workflows.
\end{abstract}

\begin{IEEEkeywords}
Concept-drift monitoring, developer tooling, LLM package hallucination, npm ecosystem, policy-as-code,
sentence embeddings, slopsquatting, software supply-chain security, static analysis, typosquatting detection.
\end{IEEEkeywords}

\section{Introduction}
\IEEEPARstart{T}{he} Node Package Manager (npm) registry hosts more than two million packages and serves
billions of downloads per month, making it the largest open-source software distribution channel in the
world~\cite{zimmermann2019}. Modern JavaScript and Node.js projects routinely depend on hundreds or
thousands of third-party packages, many pulled in transitively through a single \texttt{npm install}
command. This dependency breadth accelerates development but also creates an enormous and largely
unmonitored attack surface.

Four distinct threat classes have emerged in recent years. Supply-chain compromise occurs when an
attacker injects malicious code into a previously trusted package by compromising a maintainer's account or
pushing a weaponized version update~\cite{ohm2020,latendresse2022}. The 2018 \texttt{flatmap-stream} and
2022 \texttt{node-ipc} ``peacenotwar'' incidents illustrate how a single event can affect millions of
downstream consumers simultaneously. Typosquatting exploits the difficulty humans have in detecting small
misspellings---\texttt{lodahs} for \texttt{lodash}---by registering plausible variants of popular package
names and embedding malicious post-install scripts~\cite{vu2020}. Malicious install scripts abuse npm's
lifecycle hooks (\texttt{preinstall}/\texttt{postinstall}) to execute arbitrary code with developer-level
privileges at the moment of installation, before any application code is ever run~\cite{duan2021}. Package
hallucination (also called slopsquatting) is an emergent threat enabled by large language model (LLM) coding
assistants: when a model confidently suggests a package name that does not exist on the registry, an
attacker can register that exact name with malicious content before the developer
investigates~\cite{spracklen2025}.

Current defenses are structurally reactive. Continuous integration scanners and tools such as
\texttt{npm audit} check installed packages against vulnerability databases~\cite{zapata2018}, but by the
time a pipeline runs, a package has typically been installed locally and its scripts executed. Registry-side
malware scanning catches some malicious packages before wide adoption, yet even registry operators report
that false-positive rates for automated tools remain prohibitively high~\cite{vu2023}. Commercial tools such
as Socket.dev provide some pre-install signal but require cloud round-trips, introduce privacy concerns, and
none of them systematically addresses the LLM hallucination threat.

This paper makes the following contributions.

\begin{enumerate}
\item CIDAS, an open-source, locally-run, real-time package screening system that intercepts
\texttt{npm install} before any files are downloaded or scripts executed, operating with fail-open
semantics that guarantee developer workflow continuity.

\item A four-pillar detection architecture combining semantic project-context fit (Contextify), reputation
and typosquat analysis (Sentinel), static and AST-based script scanning with cross-version diffing (Shield),
and a two-stage, weighted, explainable Aggregator---the first system, to our knowledge, to unify all four
npm threat classes in a single pre-install gatekeeper.

\item A policy-as-code governance layer enabling teams to version-control dependency security rules
alongside their codebase, with schema validation that rejects misconfigured policies rather than silently
misapplying them.

\item A concept-drift monitor that continuously compares live \emph{per-pillar} score distributions against a
stored baseline using KL divergence, providing ongoing assurance that the scoring system remains well
calibrated without manual re-tuning.

\item Two detection-quality refinements identified and validated through a live evaluation against a real
labelled corpus rather than synthetic benchmarking alone: reputation-corroborated, affix-aware typosquat
detection, and manifest-first gating of the expensive cross-version tarball diff.

\item A comprehensive empirical evaluation across 548 automated tests, a component-wise ablation study, and
labelled corpora covering all four threat categories, demonstrating correct classification and
network-condition-dependent scan latency on commodity hardware.
\end{enumerate}

The remainder of the paper is organized as follows. Section~II surveys related work. Section~III defines the
threat model. Section~IV presents the system architecture. Section~V details each analysis pillar. Section~VI
describes supporting subsystems. Section~VII covers implementation. Section~VIII reports experimental
results. Section~IX discusses findings and limitations. Section~X concludes.

\section{Related Work}

\subsection{Supply-Chain Attacks and Package-Manager Security}
Cappos et al.~\cite{cappos2008} established the foundational security threat model for package managers,
demonstrating that man-in-the-middle and malicious-mirror attacks could compromise client machines at scale
across ten major package managers. Zimmermann et al.~\cite{zimmermann2019} conducted a landmark large-scale
analysis of npm's dependency graph, finding that fewer than 20 maintainers hold transitive influence over
more than half of all npm packages, creating concentrated systemic risk. Ohm et al.~\cite{ohm2020} developed
a comprehensive taxonomy of real-world open-source supply-chain attacks covering typosquatting, dependency
confusion, and maintainer account takeover. Latendresse et al.~\cite{latendresse2022} further characterized
maintainer account compromise as a distinct, empirically common attack vector. Duan et al.~\cite{duan2021}
empirically studied malicious package campaigns across npm, PyPI, and RubyGems, establishing that
install-time script abuse is the dominant execution vector, directly motivating the Shield pillar of CIDAS.

\subsection{Typosquatting and Name-Based Deception}
Vu et al.~\cite{vu2020} defined edit-distance heuristics for detecting typosquatted names in the Python
ecosystem, demonstrating the effectiveness of Levenshtein distance for surface-level name deception. CIDAS
originally adopted this approach directly in the Sentinel pillar; a live evaluation against a labelled corpus
(Section~VIII) subsequently showed that raw edit distance alone structurally misses affix-based squats and
false-positives on short, unrelated legitimate names, motivating the affix-canonicalization and
reputation-corroboration refinements described in Section~V-B.

\subsection{Vulnerability Propagation and Reputation Signals}
Zapata et al.~\cite{zapata2018} studied vulnerable dependency migration patterns in npm. Zahan et
al.~\cite{zahan2022} identified package-level risk indicators---download volume, maintainer count, package
age, and repository link presence---that predict supply-chain risk; these signals inform the core of CIDAS's
Sentinel reputation sub-pillar. Vu et al.~\cite{vu2021} surveyed detection feasibility across multiple attack
classes, motivating CIDAS's multi-signal aggregation strategy.

\subsection{Static Analysis and Automated Detection}
Sejfia and Sch\"{a}fer~\cite{sejfia2022} proposed Amalfi, a machine-learning system combining metadata-based,
behavioral, and static-analysis classifiers to detect malicious npm packages in seconds without execution,
closely mirroring the CIDAS Sentinel and Shield design. Ferreira et al.~\cite{ferreira2021} showed that a
lightweight per-package permission system can sandbox most npm packages at near-zero cost, informing a
future-work direction for CIDAS. Froh et al.~\cite{froh2023} explored differential static analysis for
detecting malicious updates; CIDAS's cross-version diff analyzer operationalizes this idea within the Shield
pillar, and Section~V-C describes a manifest-first gate added to make the differential check practical at
scale.

\subsection{LLM-Induced Threats}
Spracklen et al.~\cite{spracklen2025} documented the package-hallucination phenomenon at scale across
multiple LLMs, coining the term slopsquatting and quantifying the rate at which model-suggested package
names do not exist on the registry. This work is the direct motivation for CIDAS's AI-suggestion flag and its
live OSV.dev advisory check for unrecognized package names. Greshake et al.~\cite{greshake2023} described
prompt-injection patterns in LLM-integrated applications, informing CIDAS's static scanning of package
README files and install scripts.

\subsection{Semantic Analysis and Embeddings}
Reimers and Gurevych~\cite{reimers2019} introduced Sentence-BERT, the backbone of CIDAS's Contextify pillar
for project-context analysis. Feng et al.~\cite{feng2020} proposed CodeBERT as a foundational model for joint
code and natural-language understanding, providing relevant background for the optional LLM second-pass
verifier.

\subsection{Registry-Side Scanning}
Vu et al.~\cite{vu2023} studied operational malware scanning at PyPI, finding that even a 1-in-10,000
false-positive rate is prohibitively expensive for registry operators and that current automated tools exceed
this rate substantially. This empirical grounding motivates CIDAS's fail-open design, per-machine trust
lists, and emphasis on minimizing false positives. Froh et al.~\cite{froh2023} and Ladisa et
al.~\cite{ladisa2023} further systematized supply-chain attack classes and practical defensive
considerations.

\subsection{Differentiation}
CIDAS is, to our knowledge, the first system to (a) address all four threat classes in a single pre-install
gatekeeper, (b) integrate semantic project-context awareness as a first-class detection signal, (c) support
schema-validated policy-as-code governance, and (d) deploy a KL-divergence-based concept-drift monitor
operating on genuine per-pillar score history for continuous calibration assurance.

\emph{A note on scope.} CIDAS's Sentinel pillar targets name-based typosquatting and does not currently
address \emph{dependency confusion}---the distinct attack class in which a public package is registered
under a name matching a private, internal package, exploiting resolver precedence rather than name
similarity~\cite{birsan2021}. We flag this explicitly as an out-of-scope threat class in Section~IX-E rather
than implying coverage.

\section{Threat Model}
We consider an adversary whose goal is to execute arbitrary code on a developer's workstation or within a
CI/CD pipeline by inducing the installation of a malicious npm package. The attacker can control package
names, content, and registry metadata for packages they publish, but is assumed to lack control over npm's
core infrastructure, the developer's operating system, or the CIDAS daemon binary itself. We consider four
attack classes:

\textbf{A1. Typosquatting.} The attacker registers a package name visually or lexically similar to a popular
package and embeds malicious lifecycle scripts.

\textbf{A2. Malicious Install Scripts.} The attacker publishes a package containing hooks that exfiltrate
environment variables, download secondary payloads, or establish persistence.

\textbf{A3. Supply-Chain Compromise.} The attacker hijacks a legitimate maintainer account and pushes a
malicious version update to an otherwise trusted package.

\textbf{A4. LLM Package Hallucination.} An LLM coding assistant suggests a non-existent package name; the
attacker registers that name with malicious content before the developer investigates.

Out of scope: kernel-level rootkits, compromised npm infrastructure, post-install defenses,
dependency-confusion attacks against private package namespaces (see Section~II-H), and social-engineering
attacks not involving npm.

\section{System Architecture}

\subsection{Design Principles}
CIDAS is designed around five principles:

\textbf{P1. Pre-execution interception.} All analysis must complete before any package file is downloaded or
any script is executed.

\textbf{P2. Fail-open semantics.} A daemon outage must degrade to a logged warning, not a blocked workflow.

\textbf{P3. Explainability.} Every verdict must carry the specific signals that produced it so that
developers can make informed override decisions.

\textbf{P4. Privacy preservation.} Package contents and project metadata must never leave the developer's
machine.

\textbf{P5. Low friction.} Per-package scan latency must remain low on commodity hardware to ensure continued
adoption; Section~VIII-D discusses how this target is met on the cache-hit path and the network-bound reality
of the cache-miss path.

\subsection{High-Level Architecture}
Three developer-facing entry points---the npm shim, the VS Code extension, and a REST API---all communicate
with a single locally-run FastAPI daemon listening on port 7355. Inside the daemon, four coordinated pillars
(Contextify, Sentinel, Shield, Aggregator) execute concurrently and feed a unified verdict pipeline that also
consults the policy engine, trust list, scan cache, and offline-allow cache before returning a result to the
caller. Every decision is written to an append-only audit log monitored by the concept-drift subsystem.

\subsection{Request Priority Chain}
Every incoming scan request passes through the following priority chain before any pillar analysis is
invoked, short-circuiting as soon as a determination is possible: (1) policy block-list check
$\rightarrow$ forced block; (2) policy trust-list check $\rightarrow$ forced allow; (3) HMAC-verified local
trust-database lookup $\rightarrow$ allow if untampered; (4) one-hour SQLite scan cache $\rightarrow$ return
cached verdict; (5) proceed to concurrent pillar analysis. This design ensures that repeated scans of
unchanged, already-vetted dependencies impose negligible latency.

\subsection{npm Shim}
The npm shim is a thin cross-platform Node.js wrapper placed ahead of the real npm binary on the system
PATH. On every \texttt{npm install} invocation it (a) performs a SHA-256 self-integrity check to detect
tampering, (b) diffs the project's \texttt{package.json} to identify newly requested packages, (c) forwards
each new package name, resolved version, and working directory to the daemon via an authenticated REST call
to \texttt{/api/v1/scan}, (d) interprets the returned verdict, and (e) either passes through to the real npm
binary (ALLOW), prints a cautionary message and optionally prompts for confirmation (WARN), or blocks the
install with a human-readable explanation (BLOCK).

If the daemon is unreachable the shim first checks a 24-hour offline-allow cache of previously issued ALLOW
verdicts. For genuinely new packages with an unreachable daemon, installation proceeds with a clear bypass
notification written to the audit log (fail-open).

\subsection{VS Code Extension}
A \texttt{package.json} watcher diffs the file on every save and triggers asynchronous scans of newly added
entries. A terminal listener (\texttt{sentinelHook}) detects \texttt{npm install} commands typed into the
integrated terminal and flags the session as AI-suggested when applicable. A color-coded status bar polls
daemon health every 30 seconds and surfaces WARN/BLOCK dialogs with a Show Details primary action that opens
a webview panel displaying the full per-pillar score breakdown and the resolved policy file path.

\section{Analysis Pillars}

\subsection{Contextify: Project-Context Fit}
Contextify assesses whether a candidate package is semantically appropriate for the target project. A
package such as a cryptocurrency miner may carry no reputation flags and no malicious scripts, yet is
clearly anomalous in a React web-application context. Detecting such mismatch requires understanding the
project's existing dependency profile.

\subsubsection{Project Fingerprinting}
CIDAS builds a project fingerprint by scanning up to 150 JavaScript and TypeScript source files in the
target directory, extracting \texttt{require} and \texttt{import} statements, retrieving the npm registry
description of each discovered dependency, and concatenating those descriptions into a single text corpus.
This corpus is encoded into a dense vector using the all-MiniLM-L6-v2 Sentence-Transformer
model~\cite{reimers2019} (90\,MB, downloaded once, cached locally, CPU-only inference).

\subsubsection{Similarity Scoring}
The candidate package's description is independently encoded and compared against the project fingerprint
vector via cosine similarity. \texttt{compute\_score} maps this similarity to a bounded 0--25 risk
contribution: 0 when similarity $\geq 0.65$; 10 (\texttt{loosely\_related}) when similarity $\geq 0.35$; 25
(\texttt{unfamiliar\_in\_mature\_project}) when the project has more than 10 tracked dependencies and
similarity remains low; 15 (\texttt{unfamiliar\_package}) otherwise. Because a 0--25 contribution at weight
0.30 can reach at most 7.5 points---far below the 40-point WARN threshold---the Aggregator applies a
separate, weight-independent floor described in Section~V-D.

\subsection{Sentinel: Reputation and Name Analysis}

\subsubsection{Known-Incident Blocklist}
Sentinel maintains a curated list of packages with documented, severe supply-chain incidents (e.g.,
\texttt{flatmap-stream}, \texttt{event-stream}, \texttt{ua-parser-js}, \texttt{node-ipc}, \texttt{coa},
\texttt{rc}, \texttt{eslint-scope}). Any match forces an immediate score of 95 and the
\texttt{known\_supply\_chain\_incident} flag, which the Aggregator's Stage~1 gate (Section~V-D) escalates to
a forced BLOCK regardless of the weighted score.

\subsubsection{Registry Reputation}
For packages not on the blocklist, Sentinel evaluates registry signals: monthly download count (via the npm
downloads-point API), days since first publication, and repository-link presence, following the
risk-indicator framing of Zahan et al.~\cite{zahan2022}. \texttt{check\_registry\_existence} additionally
collects a per-package maintainer count into scan metadata; this signal is presently \emph{recorded but not
yet incorporated into the risk score}, and is flagged as a validated-threshold-pending refinement in
Section~IX-F.

\subsubsection{Typosquat Detection}
Sentinel applies case-normalized Levenshtein distance between the candidate name and each of approximately
53 curated high-popularity npm packages (a placeholder pending a full top-500 list). A candidate within edit
distance 2 is a \emph{raw-distance hit}. A live evaluation against a labelled corpus (Section~VIII) found
this check structurally insufficient in two directions: it cannot catch \emph{affix squats}
(\texttt{node-react} is edit distance 5 from \texttt{react}, past any workable threshold), and it
false-positives between short, unrelated, legitimate packages that happen to collide within distance~2
purely by coincidence (\texttt{vue} vs.\ \texttt{vite}, distance 2; \texttt{next} vs.\ \texttt{nuxt},
distance 1). A single global edit-distance threshold cannot resolve both failure modes simultaneously, since
tightening it worsens the first and loosening it worsens the second. CIDAS addresses this with two
complementary, decoupled refinements rather than a threshold change:

\begin{itemize}
\item \textbf{Affix canonicalization.} A candidate name is stripped of one leading squat prefix
(\texttt{node-}, \texttt{js-}) and one trailing squat suffix (\texttt{-js}, \texttt{-util(s)},
\texttt{-helper}, \texttt{-async}, \texttt{-core}, \texttt{-lib}) and compared for an \emph{exact} match
against the curated list, surfacing as \texttt{typosquat\_affix\_match}. This is a structurally distinct
signal from raw edit distance, not a threshold variant of it, and independently catches names like
\texttt{node-react} that raw distance never can.
\item \textbf{Reputation corroboration.} Neither a raw-distance nor an affix hit alone forces the maximum
score. Sentinel additionally fetches the \emph{matched target's} own download count and publication age
(reusing the same short-TTL, single-flight registry cache used elsewhere in the daemon, at no additional
network round-trip cost when the target was already fetched by a concurrent pillar) and confirms a
disparity---candidate downloads under 5\% of the target's, or a candidate younger than 30 days against a
target older than 365---before escalating to a forced score of 100 and the \texttt{typosquat\_detected}
flag. This directly suppresses the \texttt{vue}/\texttt{vite} and \texttt{next}/\texttt{nuxt} class of false
positive, since both sides of those pairs are comparably large and long-established. If the target lookup
itself fails (network error, confirmed absence), corroboration \emph{fails toward flagging}---returning
confirmed---so a corroboration-check outage never silently suppresses the pre-existing detection.
\end{itemize}

Both refinements are governed by independent per-machine admin-config flags
(\texttt{typosquat\_affix\_canonicalization}, \texttt{typosquat\_reputation\_corroboration}; both default
\texttt{true}), allowing either to be disabled instantly in the field without a code change if it misbehaves.
A candidate matching a typosquat pattern (raw or affix) that does not yet exist on the registry receives the
same maximum score (100) as an existing typosquat once corroborated---pre-registration does not reduce
severity, since an unregistered typosquat name is itself a strong signal of staged attacker intent, and both
signals are already at the scoring ceiling.

\subsubsection{AI Hallucination Detection}
When the shim or extension flags a package as \texttt{ai\_suggested} via the \texttt{sentinelHook}, Sentinel
applies stricter scrutiny: (i) a package absent from the registry sets \texttt{package\_not\_found} and
forces an automatic block; (ii) packages published under 7 days ago receive an escalated reputation
sub-score (\texttt{very\_new\_package}, +35) and those under 30 days a moderate one
(\texttt{new\_package}, +15); and (iii) a live query to the public OSV.dev vulnerability API checks for
known advisories or confirmed malware entries.

\subsection{Shield: Static and AST-Based Analysis}
Shield downloads the package tarball and analyzes its contents without executing any code.

\subsubsection{Lifecycle Script Pattern Scanning}
Shield extracts all lifecycle hooks from \texttt{package.json} and applies regex-based pattern matchers
covering seven categories: (i) environment-variable exfiltration via \texttt{process.env} combined with
outbound network calls, (ii) generic outbound network calls in an install script (\texttt{curl},
\texttt{wget}, \texttt{fetch}), (iii) Base64-encoded eval payloads (\texttt{Buffer.from}, \texttt{atob}),
(iv) child-process spawning to download or execute secondary payloads, (v) cryptocurrency-miner signature
strings, and (vi)--(vii) two classes of obfuscation (contiguous hex-escape runs and long numeric-array
literals).

\subsubsection{AST-Based Analysis}
For source files flagged as high-risk by regex pre-filtering, Shield parses JavaScript into an abstract
syntax tree (AST) using tree-sitter and performs structural analysis. AST-based analysis is more robust than
regex against obfuscation techniques that split keywords across concatenated string literals or use computed
property access.

\subsubsection{Cross-Version Diff Analysis}
For packages with a previously scanned and trusted version, Shield performs a capability diff: it compares
the new version's set of npm module imports, \texttt{process.env} accesses, and outbound network call
targets---all extracted via AST---against the immediately-preceding published version. New capabilities
trigger a \texttt{new\_network\_calls}/\texttt{new\_imports} signal and raise the Shield sub-score, enabling
detection of compromised maintainer accounts that push malicious updates to legitimate
packages~\cite{froh2023}.

This diff requires downloading and extracting \emph{two} full tarballs (current and immediately-preceding
version), which a live evaluation (Section~VIII-D) identified as the dominant contributor to per-scan
latency. CIDAS mitigates this with a \textbf{manifest-first gate}: since every published version's
\texttt{scripts} and \texttt{dependencies} fields are already present in the registry document already
fetched for other purposes, Shield compares them between the current and preceding version \emph{before}
downloading the preceding version's tarball. When both fields are byte-identical, the tarball diff is skipped
entirely---there is no declared lifecycle-script or dependency change for it to explain---while the primary,
single-tarball regex/AST scan of the current version still runs unconditionally, so absolute malicious
patterns in file content are still caught regardless of gating. The gate's accepted residual risk is narrow:
a compromised update that adds a new capability purely through code (e.g., a Node built-in
\texttt{require('dns')} call) without touching the manifest would not trigger the differential
\texttt{new\_network\_calls} escalation, though it remains subject to the unconditional absolute-pattern
scan. The gate is governed by \texttt{shield\_manifest\_gating} (default \texttt{true}) and was validated to
not regress recall on the corpus's account-takeover and protestware categories (Section~VIII-D).

\emph{Predecessor resolution past a purged version.} ``Immediately-preceding published version'' above is
resolved by \texttt{get\_previous\_version}, which now reasons about a package's full publish-time ordering
rather than requiring an exact match against the current version within an already-resolvability-filtered
list. This distinction matters specifically when npm's security team has purged a version's listing after a
confirmed compromise: the purged version's own timestamp is typically still present in the registry's publish
history even though its manifest is gone, so its position in that ordering can still be located, and the
predecessor search then walks backward from that position, skipping over any further purged versions, to the
nearest version whose manifest is still resolvable. This walk is bounded to 20 versions back, a scope decision
recorded explicitly here and in code: npm's takedown process typically purges only the one or two malicious
releases themselves, but an unbounded walk could in principle diff against a version many releases older than
intended if a very long run were purged. Section~\ref{sec:new-incidents} reports the empirical motivation and
validation of this change.

\emph{Tri-state version resolution.} Independently, \texttt{Shield.score()} no longer silently substitutes
\texttt{dist-tags.latest} when the exact requested version cannot be resolved in the registry's version map.
It instead returns an explicit \texttt{requested\_version\_unresolved} state---score 0, confidence 0, no
tarball scanned---which the Aggregator's Stage~1 gates (Section~V-D) treat as a floor at
\texttt{warn\_threshold} rather than a clean result. This mirrors the tri-state (exists / confirmed-absent /
undetermined) discipline already applied to Sentinel's registry-existence check: a scanner that cannot
examine the specific artifact it was asked to examine must say so explicitly, rather than examining something
else (here, the current clean release) and returning a verdict as though the requested version had actually
been scanned. Section~\ref{sec:new-incidents} documents the concrete case this closes.

\subsection{Aggregator: Two-Stage Weighted Score and Verdict}
The Aggregator combines the three pillar sub-scores into a composite risk score $r$ via a two-stage
structure: a deterministic Stage~1 gate supplies an optional score floor, and a weighted Stage~2 score is
always computed; the two are combined via $\max(\cdot,\cdot)$. This is organized explicitly as two named
stages---rather than a single sequential pass, as in an earlier revision of this system---because an
ablation study (Section~VIII-E) showed that Stage~1's flag-driven conditions already determine the verdict
for the large majority of a real labelled corpus, independent of pillar weighting; separating the stages
gives future signals (e.g., a learned classifier, homoglyph detection) a clean place to extend Stage~2
without diluting into a sum that Stage~1 already dominates for most records.

\textbf{Stage 2 (weighted score).}
\begin{equation}
s_2 = w_C s_C + w_S s_S + w_{Sh} s_{Sh}
\label{eq:stage2}
\end{equation}
where $w_C = 0.30$ (Contextify), $w_S = 0.35$ (Sentinel), and $w_{Sh} = 0.35$ (Shield). All sub-scores and
$s_2$ lie in $[0,100]$ with higher values indicating greater risk. Two additive modifiers apply on top of
Eq.~\eqref{eq:stage2}:

\begin{itemize}
\item \textbf{Contextify floor.} If the raw embedding cosine similarity between the candidate and the
project fingerprint falls below 0.05, a flat +20 penalty is applied and \texttt{alien\_to\_project} is
tagged---closing the gap for packages essentially unrelated to anything the project imports, which
Contextify's own bounded 0--25 contribution (Section~V-A) cannot reach on its own.
\item \textbf{Covert-dropper amplification.} Co-occurrence of a Shield \texttt{base64\_decode} flag and a
Contextify \texttt{unfamiliar\_in\_mature\_project} flag raises the score by 15 points (capped at 100),
tagging \texttt{covert\_dropper\_signals}---penalizing the pattern characteristic of covert droppers injected
into unfamiliar dependencies.
\end{itemize}

\textbf{Stage 1 (deterministic gates).} Five flag-driven conditions, evaluated independently of pillar
weight, each supply a score floor that Stage~2's result cannot fall below:

\begin{itemize}
\item \texttt{package\_not\_found} $\rightarrow$ floor at \texttt{block\_threshold} (80): a nonexistent
package is either hallucinated or malicious.
\item \texttt{known\_supply\_chain\_incident} $\rightarrow$ floor at \texttt{block\_threshold} (80).
\item \texttt{npm\_security\_placeholder\_version} $\rightarrow$ floor at \texttt{block\_threshold} (80): the
resolved version matches npm's \texttt{-security.N} placeholder convention, meaning npm's security team has
pulled this exact release as malicious or reserved and republished an inert stub in its place
(Section~\ref{sec:new-incidents}).
\item \texttt{typosquat\_detected} $\rightarrow$ floor at \texttt{warn\_threshold} (40): Sentinel's weight
alone caps its contribution at 35 points, below WARN, so a corroborated typosquat with an otherwise-quiet
Contextify/Shield signal must not silently ALLOW. WARN rather than BLOCK is used here specifically because
name-similarity, even corroborated, can still legitimately false-positive on a genuinely unrelated package;
BLOCK is reserved for the confirmed-malicious conditions above.
\item \texttt{requested\_version\_unresolved} $\rightarrow$ floor at \texttt{warn\_threshold} (40): Shield
could not resolve the specific requested version in the registry (Section~V-C) and therefore examined
nothing. This is not itself confirmed-malicious---unlike a known incident or a security placeholder---but
"the exact pinned version cannot be examined at all" is meaningful enough to warrant caution rather than a
silent pass-through on whatever Contextify and Sentinel alone conclude.
\end{itemize}

The composite score maps to verdicts as: $r \in [0,39] \Rightarrow$ ALLOW; $r \in [40,79] \Rightarrow$ WARN;
$r \in [80,100] \Rightarrow$ BLOCK.

\section{Supporting Subsystems}

\subsection{Policy-as-Code Engine}
Organizations express dependency governance rules in a \texttt{.cidas/policy.json} file co-located with
their repository. The schema (a Pydantic model, \texttt{version: 1}, \texttt{extra="forbid"}) supports:
\texttt{block\_list} (always BLOCK), \texttt{trust\_list} (bypass pillar analysis),
\texttt{min\_monthly\_downloads}, \texttt{require\_repository\_link}, \texttt{max\_sentinel\_distance} (a
per-project override of the default Levenshtein typosquat threshold), \texttt{contextify\_weight} (a
per-project override of the Contextify pillar weight, 0.0--0.5, with the remainder re-split between Sentinel
and Shield in their existing ratio), and \texttt{warn\_requires\_confirmation} (force interactive
confirmation for WARN verdicts). The policy file is validated against this schema, which rejects
unrecognized fields with a hard error rather than silently ignoring them, preventing governance bugs caused
by typos.

\subsection{HMAC-Protected Trust List}
Trust-list entries are stored in SQLite with HMAC-SHA256 integrity tags derived from a locally generated
machine secret. On every lookup the stored tag is re-verified; a mismatch triggers a CRITICAL audit-log event
and falls through to a full pillar scan.

\subsection{Two-Tier Cache}
(1) A one-hour scan cache stores complete verdicts and pillar sub-scores keyed by (name, version) pairs,
bounding staleness for rapidly evolving threat intelligence. (2) A 24-hour offline-allow cache stores ALLOW
verdicts issued while the daemon was reachable, enabling previously approved packages to install without
interruption during daemon outages. An explicit \texttt{/api/v1/cache/invalidate} endpoint forces re-analysis
of a specific package on the next request; a separate \texttt{DELETE /api/v1/cache} endpoint purges expired
entries.

Beneath the scan-level cache, a shorter-lived (10-second), single-flight cache deduplicates the underlying
npm registry document fetch itself: a single scan of an existing package independently needs the same
package's full registry document from up to five call sites across the daemon (Contextify's description
lookup, Sentinel's existence/reputation check, Shield's tarball resolution, the disk-footprint estimator, and
the diff analyzer's version-history lookup). Without deduplication this produced measurably redundant,
concurrent requests to the same registry endpoint within a single scan (Section~VIII-D); the single-flight
cache collapses concurrent and near-concurrent requests for the same package name into one underlying fetch,
and does not cache a failed fetch, so a transient registry outage cannot poison the result for other callers
within the TTL window.

\subsection{Append-Only Audit Log}
Every scan decision---verdict, per-pillar scores, active flags, policy file path, user overrides, and bypass
events---is appended as a newline-delimited JSON (JSONL) record to a rotating log file (10\,MB per file,
three-file rotation, four files total including the active log). Each record now persists the individual
Contextify, Sentinel, and Shield scores alongside the composite verdict, which is what makes genuine
per-pillar concept-drift monitoring possible (Section~VI-E). The log is queryable via a CLI and a REST
endpoint supporting filters on verdict, package name, and time range.

\subsection{Concept-Drift Monitor}
On each health-check poll the monitor: (1) reads the most recent 100 audit-log records, each of which now
carries all three per-pillar scores; (2) computes 10-bin per-pillar score histograms; (3) computes KL
divergence $D_{KL}(P\|Q)$ between the live histogram $P$ and a stored baseline histogram $Q$ for each pillar;
(4) returns \texttt{drift: "warn"} if the pillar-averaged divergence exceeds 0.15, or \texttt{drift: "alert"}
if it exceeds 0.30.

Because no production deployment history yet exists to build a baseline from, the baseline is instead seeded
synthetically: an offline build step assigns each record in the labelled evaluation corpus a per-pillar score
sampled from an empirically-informed range keyed to its ground-truth label (e.g., malicious packages sampled
high on Sentinel and Shield; benign packages sampled low across all three), and builds the baseline histogram
from that synthetic distribution rather than from real traffic. This is an explicit, documented placeholder:
Section~IX-F recommends replacing it with a baseline built from real deployment history once one accumulates.
With the per-pillar score persistence described above now in place, a live health check against real scan
traffic correctly reports \texttt{status: "alert"} with a non-trivial KL divergence against this synthetic
baseline (Section~VIII-F)---confirming the mechanism is now genuinely operative, whereas in an earlier
revision of the audit-log schema (which persisted only the composite score) per-pillar drift detection was
structurally impossible and the monitor always reported \texttt{insufficient\_data} against real traffic.

\section{Implementation}

\subsection{Technology Stack}
Table~\ref{tab:stack} summarizes the primary technologies used in each CIDAS component.

\begin{table}[t]
\caption{CIDAS Technology Stack}
\label{tab:stack}
\centering
\begin{tabular}{@{}p{0.22\linewidth}p{0.68\linewidth}@{}}
\toprule
\textbf{Component} & \textbf{Key Technologies} \\
\midrule
Daemon & Python 3.10+, FastAPI, Pydantic v2, SQLite, all-MiniLM-L6-v2 (HuggingFace sentence-transformers), tree-sitter, httpx \\
VS Code Ext. & TypeScript 5, VS Code Extension API (v1.89+), vitest \\
npm Shim & Node.js 18+, cross-platform bash / PowerShell 5.1+ \\
Testing & pytest / pytest-cov (daemon), vitest (extension) \\
\bottomrule
\end{tabular}
\end{table}

\subsection{Daemon and REST API}
The daemon exposes nine endpoints under \texttt{/api/v1/} (Table~\ref{tab:api}, Appendix~A). Pillar functions
are dispatched concurrently via \texttt{asyncio.gather}, maximizing throughput within the latency budget.

\subsection{Cross-Platform Shim}
The shim (\texttt{npm-shim.js}) supports macOS, Linux, WSL~2, and native Windows. Platform-specific
installation scripts discover the real npm binary, rename it, install the shim in its place, and compute the
SHA-256 reference hash used for self-verification.

\subsection{Optional LLM Second-Pass Verifier}
For ambiguous README or install-script text, CIDAS optionally invokes a locally running Ollama-hosted LLM
(default \texttt{phi3:mini}). The LLM score is blended at $0.60 \times \text{LLM} + 0.40 \times \text{regex}$.
The regex pass is high-precision but low-recall---it reliably catches canonical phrasings such as ``ignore
previous instructions'' but misses paraphrases---while the LLM exhibits the opposite profile. Weighting the
LLM higher recovers the recall the regex pass misses, while the regex's 40\% floor still anchors the blended
score against a purely hallucinated LLM judgment. All inference is local; no package content leaves the
developer's machine.

\section{Evaluation}

\subsection{Experimental Setup}
All experiments were conducted on a commodity developer workstation without a GPU, against the daemon
running locally and a live connection to the public npm registry and the npm downloads API---not a mocked
or air-gapped environment. The all-MiniLM-L6-v2 model was loaded on first run and cached locally. The
548-test suite (449 daemon, pytest; 99 extension, vitest) was executed as part of continuous verification
alongside the corpus evaluation below. Because latency in this configuration is dominated by live
npm-registry and tarball-download round-trips rather than daemon compute, absolute latency numbers depend on
registry responsiveness and network contention at scan time and should be read as characteristic ranges, not
fixed guarantees.

\subsection{Labelled Evaluation Corpus}
We constructed four disjoint corpora, generated and version-controlled by a corpus-build script
(\emph{N} = 179 total): (1) \textbf{Benign} (50 packages): top-500 npm packages by download count; (2)
\textbf{Typosquat} (50 packages): synthetic mutations (deletion, insertion, qwerty-adjacent substitution,
homoglyph, prefix, suffix) of ten high-popularity target packages, each record carrying an
\texttt{npm\_registered} flag recording whether the generated squat is actually published; (3)
\textbf{Hallucinated} (23 packages): synthetically generated using suffix templates commonly produced by LLM
code assistants (\texttt{<lib>-helper}, \texttt{<lib>-async}, \texttt{<lib>-utils}), each verified absent
from the registry at corpus-build time; (4) \textbf{Malicious} (56 packages): 39 real disclosed npm
supply-chain incidents from the corpus's original construction (event-stream/flatmap-stream, eslint-scope,
the Oct--Nov 2021 ua-parser-js/coa/rc cluster, node-ipc/peacenotwar, colors/faker, and the May-2018
getcookies credential-theft family) plus the known-incident blocklist, extended with \textbf{17 additional
records} covering real 2025--2026 supply-chain incidents added in a later pass of this evaluation: the
September 2025 Shai-Hulud stolen-token worm cluster (\texttt{@ctrl/tinycolor}, \texttt{@ctrl/deluge},
\texttt{angulartics2}, \texttt{koa2-swagger-ui}, \texttt{json-rules-engine-simplified}; 6 records), the
2026 TanStack CI/OIDC pipeline-compromise chain (\texttt{@tanstack/react-router},
\texttt{@tanstack/router-core}, \texttt{@tanstack/arktype-adapter}, \texttt{@tanstack/history}; 5 records),
the September 2025 compromised-maintainer-account incidents affecting \texttt{chalk} and \texttt{debug}, the
2026 \texttt{axios} maintainer-credential-theft incident (2 versions), and \texttt{plain-crypto-js} (2
versions, an npm-security-pulled malicious dependency injected via the March 2026 axios compromise). This
17-record extension is the direct subject of the version-resolution investigation in
Section~\ref{sec:new-incidents}.

\subsection{Detection Performance}
Table~\ref{tab:detection} reports, per ground-truth category, the verdict breakdown (ALLOW/WARN/BLOCK counts)
and the derived recall (or false-positive rate, for the negative-class benign row) and F1, measured with a
live run of the daemon against the full 179-record corpus. Counts are taken directly from the evaluation
harness's ground-truth-keyed confusion matrix, so every row's counts sum exactly to that category's N; we
deliberately avoid the harness's secondary \texttt{attack\_type} breakdown for this table, since several
attack-type tags (notably \texttt{typosquat} and \texttt{none}) are shared across corpora and would otherwise
mix records from different ground-truth categories into one number.

\begin{table}[t]
\caption{Detection Performance on the Labelled Corpus (N = 179)}
\label{tab:detection}
\centering
\begin{tabular}{@{}lccccc@{}}
\toprule
\textbf{Category} & \textbf{ALLOW} & \textbf{WARN} & \textbf{BLOCK} & \textbf{R / FPR} & \textbf{F1} \\
\midrule
Malicious (N=56)    & 12 & 20  & 24 & 0.79          & 0.88 \\
Typosquat (N=50)    & 0  & 30 & 20 & 1.00          & 1.00 \\
Hallucinated (N=23) & 5  & 1  & 17 & 0.78          & 0.88 \\
Benign (N=50)       & 50 & 0 & 0 & 0.00\textsuperscript{a} & ---\textsuperscript{b} \\
\midrule
\textbf{Overall}    & \multicolumn{3}{c}{---} & \textbf{0.87}\textsuperscript{c} & \textbf{0.93}\textsuperscript{d} \\
\bottomrule
\end{tabular}
\\[2pt]
\raggedright\footnotesize
\textsuperscript{a}Reported as false-positive rate (WARN+BLOCK $/$ 50), not recall---benign is a
negative-class-only corpus, so it has no recall to report. All 50 benign records resolved to a verdict in
this run (zero scan errors).
\textsuperscript{b}Precision and F1 are undefined for an all-negative corpus (there are no true positives to
define them against); FPR is the appropriate single-number summary here, consistent with why this row is
excluded from the macro-F1 computed below.
\textsuperscript{c}Recall (aggregate F1's recall component) over all 129 actually-positive records
(malicious + typosquat + hallucinated) combined.
\textsuperscript{d}\emph{Aggregate F1}: a single pooled, micro-averaged F1 over all 179 records (WARN/BLOCK
as the positive prediction, ALLOW as negative), \emph{not} the mean of the four rows above---see the
macro-F1 discussion immediately below for that alternative statistic.
\end{table}

\emph{Aggregate F1 vs.\ macro-F1.} The headline 0.93 figure reported above and in the Abstract is a pooled,
micro-averaged F1 across all 179 records treating WARN/BLOCK as the positive prediction---it is influenced by
corpus size (the 50-record typosquat corpus, at perfect recall, and the 56-record malicious corpus outweigh
the smaller 23-record hallucinated corpus). A category-balanced \emph{macro-F1}, computed as the unweighted
mean of the three F1-bearing rows in Table~\ref{tab:detection} (malicious 0.88, typosquat 1.00, hallucinated
0.88; benign is excluded per footnote b, since F1 is undefined for an all-negative corpus), is
\textbf{0.92}---close to the pooled figure here, but the two statistics answer different questions and should
not be conflated; we report both rather than picking whichever is larger.

A methodological caveat on the perfect precision (1.00) and zero false-positive rate reported above: the
benign corpus (Section~VIII-B) consists of the top-500 npm packages by download count---i.e., established,
high-reputation packages that are, by construction, comfortably distant from every pillar's risk signals
(recent, low-download, unfamiliar-to-the-project, or structurally similar to a known-bad pattern). A
zero-FPR result on this specific corpus is evidence that CIDAS does not spuriously flag well-established,
popular packages; it is not evidence of a zero false-positive rate on the long tail of low-download,
new, or idiosyncratic-but-legitimate packages, which this corpus does not sample and which is a materially
harder population for any reputation- or novelty-sensitive signal to avoid over-flagging.

Two refinements to Sentinel's typosquat detection (Section~V-B) were made in direct response to this
evaluation and are reflected in Table~\ref{tab:detection}'s typosquat row. To isolate their effect cleanly,
we reran the typosquat corpus alone with both refinements disabled via their admin-config kill-switches
(Section~V-B), holding every other part of the system---including the earlier WARN-floor aggregator
fix---fixed: with affix canonicalization and reputation corroboration off, raw Levenshtein distance alone
correctly escalated 40 of 50 records (recall 0.80, confusion matrix ALLOW=10/WARN=20/BLOCK=20), missing every
prefix/suffix-mutation record in the corpus (the structural blind spot described in Section~V-B) and
false-positively forcing a BLOCK-level score on legitimate, unrelated packages such as \texttt{vue}. With both
refinements enabled (the production default), the same 50-record corpus reaches \textbf{recall 1.00}
(confusion matrix ALLOW=0/WARN=30/BLOCK=20)---every affix squat that raw distance missed is now caught by
affix canonicalization, and reputation corroboration confirms each hit without introducing a new false
positive on this corpus. The benign corpus's false-positive rate over the same period fell from 0.15 (an
earlier, less-refined measurement, prior to the WARN-floor and corroboration fixes) to the 0.02 reported in
Table~\ref{tab:detection}---consistent with corroboration suppressing the short-name-coincidence class of
false positive described in Section~V-B, rather than trading recall for precision.

We note one earlier imprecision in reporting this comparison during development of this system: an
intermediate measurement conflated the evaluation harness's \texttt{attack\_type}-keyed breakdown (which
mixes a small number of \texttt{malicious.jsonl} records tagged \texttt{attack\_type: "typosquat"} into the
same bucket as the 50-record typosquat corpus, inflating the denominator to 57) with the ground-truth-keyed
corpus count used here. The corrected, single-corpus comparison above (0.80 $\rightarrow$ 1.00) supersedes
that intermediate figure; we surface this explicitly rather than silently revising the number, consistent
with this paper's broader practice of reporting evaluation-process corrections alongside results.

Table~\ref{tab:detection}'s malicious row (recall 0.79) aggregates considerable sub-type variance, which
Table~\ref{tab:malicious-subtypes} breaks out by \texttt{attack\_type}.

\begin{table}[t]
\caption{Malicious Sub-Type Breakdown\textsuperscript{a}}
\label{tab:malicious-subtypes}
\centering
\begin{tabular}{@{}lcccc@{}}
\toprule
\textbf{Attack Sub-Type} & \textbf{N} & \textbf{P} & \textbf{R} & \textbf{F1} \\
\midrule
Dependency hijack       & 1  & 1.00 & 1.00 & 1.00 \\
Install-script execution & 2  & 1.00 & 1.00 & 1.00 \\
Account takeover        & 20 & 1.00 & 1.00 & 1.00 \\
Worm propagation\textsuperscript{b} & 6 & 1.00 & 1.00 & 1.00 \\
CI/OIDC pipeline compromise\textsuperscript{b} & 5 & 1.00 & 1.00 & 1.00 \\
Protestware             & 7  & 1.00 & 0.71 & 0.83 \\
Credential theft        & 8  & 1.00 & 0.38 & 0.55 \\
\bottomrule
\end{tabular}
\\[2pt]
\raggedright\footnotesize\textsuperscript{a}These seven \texttt{attack\_type} values are safe to source from
the harness's attack-type breakdown because each is exclusive to \texttt{malicious.jsonl} (unlike
\texttt{typosquat}/\texttt{none}, discussed above). They cover 49 of the 56 malicious records; the remaining
7 carry \texttt{attack\_type: "typosquat"} (hijacked-account releases published under a name resembling the
original package) and are included in Table~\ref{tab:detection}'s aggregate malicious row but omitted here to
avoid double-counting against the typosquat corpus.
\textsuperscript{b}These two sub-types comprise the 11 of the 17 newly-added 2025--2026 incident records
discussed in Section~\ref{sec:new-incidents} (the remaining 6---\texttt{chalk}, \texttt{debug}, two
\texttt{axios} versions, and two \texttt{plain-crypto-js} versions---are tagged \texttt{account\_takeover}
and included in that row above).
\end{table}

Credential theft is now the single lowest-recall sub-type. Protestware recall rose substantially (0.43
$\rightarrow$ 0.71) in this pass for a specific, attributable reason described in
Section~\ref{sec:new-incidents} (two of the seven protestware records, both pinned to a since-purged version,
newly benefited from the same fix that resolves the 17-record finding below), not from a change to Shield's
behavioral-pattern coverage. Credential theft's remaining gap, like the residual protestware misses
(\texttt{peacenotwar}, \texttt{faker}), reflects packages whose distinguishing signal is behavioral (a
lifecycle script that behaves destructively or exfiltrates credentials only under specific runtime
conditions) rather than purely name- or static-pattern-based, and is discussed as a limitation in
Section~IX-E; Shield's static/AST analysis and Sentinel's reputation checks are, by design, not expected to
catch a behaviorally-triggered payload that carries no distinguishing static signature.

\subsection{Malicious Corpus Extension: Real 2025--2026 Incidents}
\label{sec:new-incidents}
The 17-record malicious-corpus extension described in Section~VIII-B (Section~\ref{sec:new-incidents} is this
section) was initially misclassified by every pillar (recall 0.00 on this subset in an earlier pass of this
evaluation), and that failure was initially attributed to a structural blind spot: npm's security team
purges a compromised release's version listing once a supply-chain incident is confirmed, and the
cross-version capability diff (Section~V-C) was assumed to have no historical predecessor left to compare
against. Direct investigation of this claim against the live registry showed it to be false. All 15 affected
records had a perfectly valid, resolvable predecessor version sitting immediately adjacent to the purged one
(e.g., \texttt{axios@1.14.1}'s predecessor \texttt{axios@1.14.0} was, and remains, fully resolvable). The
actual root cause was two independent, previously undiagnosed bugs in the diff and version-resolution
pipeline, both since fixed and documented as an architectural change in Section~V-C:

\begin{enumerate}
\item \texttt{get\_previous\_version}'s lookup required the requested (purged) version to itself appear in
an already-resolvable-filtered version history before searching for its predecessor---since a purged version
by definition never appears in that filtered list, the lookup returned no predecessor unconditionally,
regardless of how much valid history existed adjacent to the gap.
\item \texttt{Shield.score()} silently substituted \texttt{dist-tags.latest} whenever the requested version
did not resolve, and proceeded to scan that unrelated, current, clean release---returning a normal-looking
verdict with no indication that a substitution had occurred. Verified directly: requests for
\texttt{axios@1.14.1}, \texttt{chalk@5.6.1}, \texttt{debug@4.4.2}, and \texttt{@ctrl/tinycolor@4.1.1} were
each found to have actually scanned a newer, unrelated, clean tarball (\texttt{axios-1.18.1.tgz},
\texttt{chalk-5.6.2.tgz}, \texttt{debug-4.4.3.tgz}, \texttt{tinycolor-4.2.0.tgz} respectively) rather than the
malicious version pinned in the corpus record.
\end{enumerate}

With both bugs fixed, a further check established a genuine, narrower limitation that the original
diagnosis had not identified: for all 15 affected records, the malicious tarball's raw bytes---not merely
its listing in the registry's version metadata---return HTTP~404 when fetched directly from npm's CDN. npm's
takedown process for a confirmed-malicious release deletes the artifact itself, not only its manifest entry.
This means no genuine AST capability diff against the actual malicious payload is possible for any of these
15 records, regardless of the two bugs above being fixed---a materially narrower and more precise limitation
than ``purged packages cannot be diffed'' in general: it is specifically packages whose tarball bytes were
deleted at the CDN level, which is npm's standard, permanent takedown behavior for confirmed-malicious
releases, that cannot be diffed. Detection for these 15 records instead comes from a new, explicit Stage~1
gate (Section~V-D): when Shield cannot resolve the requested version at all, it now returns
\texttt{requested\_version\_unresolved} and the Aggregator floors the verdict at \texttt{warn\_threshold}
(Section~V-C). This is an honest ``could not examine this artifact'' signal, not a working capability diff;
the 17-record recall figure below reflects this gate, not a successful diff of the malicious content.

Re-running the full corpus with both fixes applied yields \textbf{recall 1.0000 (17/17)} on the full extended
subset (15 originally-diagnosed records plus 2 \texttt{plain-crypto-js} records, discussed below, which share
the same version-unresolvable root cause). Applying the identical, general-purpose fix also recovered three
records in the \emph{original} 39-record malicious corpus that had the same undiagnosed problem
(\texttt{colors@1.4.1}, \texttt{colors@1.4.2}, and \texttt{eslint-config-eslint}, none of which are on the
known-incident blocklist and so were not caught by that separate mechanism), raising the original-39 subset's
recall from the previously-reported \textbf{0.62} to \textbf{0.6923 (27/39)}---a correction to a
previously-published figure, not a new measurement on new data, and reported here explicitly per this
paper's practice of surfacing evaluation-process corrections rather than silently revising them. We verified
this recall improvement is not a side effect that regressed any other corpus: the typosquat, hallucinated,
and benign confusion matrices are unchanged from the pre-fix run.

A related, separately-diagnosed and separately-fixed issue surfaced during this same investigation:
\texttt{plain-crypto-js} (both corpus versions) resolves via HTTP~200 to a tarball literally named following
npm's \texttt{-security.N} placeholder-version convention---the standard signal that npm's security team has
pulled a version and republished an inert stub in its place. Sentinel did not previously check for this
convention and scored the placeholder as an ordinary low-reputation package (ALLOW). A new Stage~1 gate
(Section~V-D) now detects this pattern directly against the resolved \texttt{dist-tags.latest} value (not the
originally-requested version string, since npm frequently wipes the entire version history down to the
single placeholder) and floors the verdict at \texttt{block\_threshold}. This is reported as a distinct
finding from the \texttt{requested\_version\_unresolved} gate above, since it addresses a different failure
mode (a resolvable-but-inert placeholder release, rather than an unresolvable one).

A separate, unrelated bug was identified and fixed during this evaluation: two entirely legitimate, popular
packages (\texttt{redux-thunk}, \texttt{nodemailer}) were observed force-BLOCKed with a
\texttt{package\_not\_found} flag in an early run. The root cause was a transient registry network failure
(timeout or dropped connection) being indistinguishable from a confirmed HTTP~404 in
\texttt{check\_registry\_existence}. This is now structurally solved: the registry-metadata function's return
type was changed from a binary exists/absent result to a tri-state \texttt{RegistryLookup} value
(\texttt{EXISTS} / \texttt{CONFIRMED\_ABSENT} / \texttt{UNDETERMINED}), threaded through every call site that
previously conflated the two. A confirmed-absent result now requires an actual HTTP~404; a timeout, a
transport error after retries are exhausted, or any other non-404 HTTP status yields \texttt{UNDETERMINED}
instead, which Sentinel routes to the existing fail-open path (a distinct, non-blocking
\texttt{registry\_lookup\_undetermined} flag) rather than the \texttt{package\_not\_found} Stage~1 gate
(Section~V-D). \texttt{check\_registry\_existence} additionally performs one confirmatory re-fetch,
bypassing the cache, before accepting a \texttt{CONFIRMED\_ABSENT} result, subsuming the earlier
cache-eviction stopgap. A related, independently discovered issue surfaced while validating this fix under
live load: npm's downloads API rate-limits aggressively under sustained request volume, which was
degrading reputation-corroboration accuracy during full-corpus runs. This was closed with two additions: a
result cache for \texttt{get\_download\_count} (5-minute TTL, since the same small set of popular target
packages is looked up repeatedly across a scan session) and a shared token-bucket rate limiter around every
outbound registry call, pacing requests proactively rather than only retrying reactively after a
429 response.

\subsection{Scan Latency}
Table~\ref{tab:latency} reports end-to-end per-package scan latency (shim call to verdict return), measured
across live runs against the full corpus.

\begin{table}[t]
\caption{Per-Package Scan Latency Under Live Registry Conditions}
\label{tab:latency}
\centering
\begin{tabular}{@{}lcc@{}}
\toprule
\textbf{Scenario} & \textbf{Median} & \textbf{P95} \\
\midrule
Cache hit (SQLite)                    & $<$50\,ms         & $<$50\,ms \\
Cache miss, warm registry conditions  & 2.7\,s            & 6.4\,s \\
Cache miss, cold cache, concurrency=4 & 19.7\,s           & 43.6\,s \\
\bottomrule
\end{tabular}
\end{table}

The cache-hit path satisfies the sub-two-second design target (P5) comfortably; the cache-miss path,
dominated by live registry round-trips and tarball downloads rather than daemon compute, does not, and its
absolute value varies substantially with registry contention (a factor of roughly 7$\times$ was observed
between the best- and worst-case runs in this evaluation, with no code change between them). Two latency
root causes were identified and addressed during this evaluation:

\begin{enumerate}
\item \textbf{Redundant registry fetches.} A single scan of an existing package was independently
re-fetching the same package's full registry document from up to five call sites with zero sharing between
them, confirmed via daemon logs showing more than nine identical requests for one package within a single
scan. The single-flight cache described in Section~VI-C reduced an initial, uncached measurement of 38.2\,s
median (80.6\,s P95) to the cache-partially-warm figures in Table~\ref{tab:latency}.
\item \textbf{Mandatory double-tarball diff.} Shield's cross-version diff (Section~V-C) downloads two full
tarballs per scan of an existing package regardless of whether the update is behaviorally meaningful. The
manifest-first gate described in Section~V-C skips this when the manifest is unchanged; a guard check
confirmed this did not regress recall on the corpus's account-takeover (0.93) or protestware (0.71)
categories, the two categories most dependent on the differential signal.
\end{enumerate}

\subsection{Ablation Study}
\label{sec:ablation}
We evaluated eight pillar-weight configurations: no pillars (baseline, $w_C{=}0/w_S{=}0/w_{Sh}{=}0$), each
pillar in isolation, each pairwise combination, and the production configuration
($w_C{=}0.30/w_S{=}0.35/w_{Sh}{=}0.35$). This design directly measures each pillar's standalone and
pairwise contribution, rather than sweeping a single weight in isolation. Table~\ref{tab:ablation} reports
results from a live run against the full corpus using the current system (i.e., reflecting the Section~V-B
and V-C refinements).

\begin{table}[t]
\caption{Pillar-Weight Ablation (N = 179)}
\label{tab:ablation}
\centering
\begin{tabular}{@{}lcccc@{}}
\toprule
\textbf{Configuration} & \textbf{P} & \textbf{R} & \textbf{F1} & \textbf{FPR} \\
\midrule
No pillars              & 0.00 & 0.00 & 0.00 & 0.00 \\
Contextify only         & 1.00 & 0.87 & 0.93 & 0.00 \\
Sentinel only           & 1.00 & 0.87 & 0.93 & 0.00 \\
Shield only             & 0.99 & 0.87 & 0.93 & 0.02 \\
Contextify + Sentinel   & 1.00 & 0.87 & 0.93 & 0.00 \\
Contextify + Shield     & 1.00 & 0.87 & 0.93 & 0.00 \\
Sentinel + Shield       & 1.00 & 0.87 & 0.93 & 0.00 \\
\textbf{All three (production)} & \textbf{1.00} & \textbf{0.87} & \textbf{0.93} & \textbf{0.00} \\
\bottomrule
\end{tabular}
\end{table}

Six of the seven non-baseline configurations---including the production weighting---now score identically at
four significant figures (P=1.0000, R=0.8682, F1=0.9295, FPR=0.0000); \emph{Shield only} is the sole
measurably-different row (F1=0.9256, FPR=0.0200), a gap too small to register at the table's two-decimal
precision but still real (Shield alone lacks Sentinel's typosquat and registry-absence Stage~1 gates). An
earlier pass of this ablation had additionally shown \emph{All three} scoring measurably \emph{worse} than
\emph{Sentinel only} (F1 0.83 vs.\ 0.84 at the time); that discrepancy was root-caused to a synchronization
bug in this evaluation harness's own local reimplementation of the Aggregator's Stage~1 gates (\S\ref{sec:ablation}
recomputes verdicts locally from stored pillar scores rather than calling the Aggregator directly, and its
hand-copied gate list had not been updated with a newly added gate), not to a genuine weighting effect; fixing
the harness's gate list to match the Aggregator resolved the discrepancy entirely, and it is not present in
Table~\ref{tab:ablation} above. This is not a null result but a specific, actionable finding about the current
aggregation design: five flag-driven Stage~1 conditions in the Aggregator (Section~V-D)---a corroborated
typosquat, a registry-confirmed-absent package, a known supply-chain incident, an npm security-placeholder
version, and an unresolvable-requested-version signal---fire independent of pillar weight and determine the
verdict for the large majority of this corpus (all of the typosquat and hallucinated sub-corpora, and most of
the malicious sub-corpus). Pillar-weight configuration is decisive only for the residual minority of records
that trigger none of these conditions,
consistent with Shield's absolute pattern-matching signal (independent of the other two pillars' flags) being
the one axis along which a configuration change measurably moves the outcome. We therefore do not claim that
the production weighting $(0.30/0.35/0.35)$ is uniquely optimal on this corpus---no configuration in
Table~\ref{tab:ablation} is meaningfully worse on F1---but retain it as the production default because it keeps
Contextify's semantic signal active for the residual population where weighting does matter, consistent with
the original design rationale in Section~V-D.

\subsection{Concept-Drift Monitor Validation}
With per-pillar score persistence in place (Section~VI-E), a live health-check poll against real scan
traffic from this evaluation correctly reported \texttt{status: "alert"} with a pillar-averaged KL divergence
in the 1.7--2.8 range against the synthetic seed baseline---confirming the monitor is now genuinely
operative against live data, a capability that was structurally absent in an earlier revision of the system
(Section~VI-E). A quantitative validation against a controlled, synthetically-shifted distribution (matching
the original design intent of injecting known-shifted scores and confirming both \texttt{alert} onset and
\texttt{ok} recovery) is identified as a remaining verification step in Section~IX-F, since the corpus-driven
traffic used here reflects a genuine distribution shift from the synthetic baseline rather than a controlled
injection.

\subsection{Test Suite Coverage}
The 548-test suite achieves 95\% statement coverage on the daemon (2,262 statements, 118 missed) and 81\%
on the VS Code extension (the untested \texttt{extension.ts} activation entry point contributes 0\% and pulls
the average down; every other tested module measures 87--100\%). Table~\ref{tab:tests} summarizes suite
composition.

\begin{table}[t]
\caption{Automated Test Suite Summary}
\label{tab:tests}
\centering
\begin{tabular}{@{}lcc@{}}
\toprule
\textbf{Suite} & \textbf{Tests} & \textbf{Coverage} \\
\midrule
Daemon (pytest, 20 modules)        & 449 & 95\% \\
VS Code extension (vitest, 5 modules) & 99  & 81\% \\
\midrule
\textbf{Total}                     & \textbf{548} & --- \\
\bottomrule
\end{tabular}
\end{table}

\section{Discussion}

\subsection{Achievement of Design Principles}
Pre-execution interception (P1) was confirmed by observing that no npm file download occurs before a BLOCK
verdict is returned, and was explicitly weighed against latency when designing the manifest-first gate
(Section~V-C), which was deliberately scoped to skip only the differential tarball fetch, never the primary
scan of the version actually being installed. Fail-open semantics (P2) were validated by bringing the daemon
offline and confirming that the shim consulted the offline cache and proceeded for known packages, issuing a
logged bypass for unknown ones. Explainability (P3) was confirmed through per-pillar scores and specific flag
codes in every WARN dialog and REST response, including the new \texttt{typosquat\_affix\_match} and
\texttt{reputation\_disparity\_confirmed} flags. Privacy preservation (P4) is ensured architecturally: the
daemon makes outbound calls only to the npm registry, OSV.dev, and HuggingFace (one-time model download).
Low friction (P5) is met on the cache-hit path; Section~VIII-D documents that the cache-miss path does not
meet it unconditionally, and characterizes the network-dependent reasons why.

\subsection{Fail-Open vs.\ Security Trade-off}
The fail-open philosophy introduces a deliberate tension with maximum enforcement. Our resolution---logging
every bypass event to the tamper-evident audit log---preserves visibility without blocking developer
workflows. This trade-off is appropriate for a developer-workstation tool; a complementary hard-fail CI gate
would provide defense-in-depth. The manifest-first gate (Section~V-C) introduces a related, narrower
trade-off deliberately accepted in this design: it exchanges a small amount of differential-detection
coverage (a manifest-invisible capability change) for a substantial latency reduction on the common case,
while leaving the absolute, non-differential detection signal fully intact regardless of gating.

\subsection{False Positives and Corroboration}
The dominant false-positive driver identified in this evaluation was not a reputation-threshold issue but a
structural one: a single global edit-distance threshold cannot simultaneously catch affix squats and avoid
short-name collisions (Section~V-B). Reputation corroboration resolved this by decoupling the two failure
modes into independent, complementary signals rather than retuning the shared threshold, and is consistent
with the observation (Section~VIII-C) that recall and the false-positive rate improved \emph{together}
rather than trading off against each other---evidence that the two original failure modes were largely
independent defects, not opposite ends of one dial.

\subsection{Observations on Multi-Pillar Design}
The combination of independent, complementary signals proved more robust than any single heuristic; this
observation is consistent with findings on multi-signal aggregation in the malware detection
literature~\cite{sejfia2022}. The ablation study (Section~VIII-E), however, surfaces a more specific finding:
for this corpus, five flag-driven Stage~1 conditions in the Aggregator (Section~V-D) already determine the
verdict for the large majority of records, and pillar-weight configuration is decisive only for the minority
of packages that trigger none of them. This is the direct motivation for formalizing the two-stage Aggregator
structure (Section~V-D) rather than treating it as an implementation detail: it makes explicit which part of
the system decides most verdicts today, and gives future weighted or learned signals (Section~IX-F) a clean
place to matter for the residual population without being diluted by conditions that already dominate.

\subsection{Limitations}
\begin{enumerate}
\item \textbf{No field deployment.} Conclusions about long-term false-positive tolerance and developer
adoption are based on a single labelled corpus and live-daemon evaluation, not longitudinal production data.
\item \textbf{Dependency confusion is out of scope.} CIDAS's Sentinel pillar targets name-based
typosquatting; it does not address dependency confusion (a public package registered under a name matching a
private internal package), a distinct and empirically impactful attack class~\cite{birsan2021} that exploits
resolver precedence rather than name similarity, and would require a different detection mechanism (private
registry/namespace awareness) entirely.
\item \textbf{Homoglyph/Unicode-confusable squatting is partially addressed, with a deliberately narrow
scope.} Sentinel now normalizes a candidate name's Cyrillic and Greek confusable characters (e.g.\ a
Cyrillic look-alike of the Latin letter ``a'' mapping to ``a'' itself) to their ASCII-skeleton equivalent
before edit-distance and affix-canonicalization comparison, catching visually-identical substitutions in
those two scripts specifically. This is an explicit, documented scope decision rather than full
Unicode-confusable coverage: the mapping table is small and hand-curated (a first-pass covering the character
blocks historically observed in registered-package-name attacks) rather than a complete import of the Unicode
confusables data set, so other scripts (e.g.\ Armenian, Cherokee) remain unaddressed and are named as future
work in Section~IX-F rather than silently uncovered.
\item \textbf{Residual typosquat recall gap.} Reputation corroboration is intentionally conservative: a
squat that has already accrued meaningful downloads or age before detection will not be corroborated as
disparate, and is therefore not force-escalated by this signal alone (Section~VIII-C).
\item \textbf{Static analysis evasion.} Novel or heavily obfuscated payloads may evade the Shield regex and
tree-sitter pattern set; the rule set requires continuous expansion.
\item \textbf{No capability diff is possible against a CDN-deleted malicious artifact.} For a package version
that npm's security team has taken down, the artifact's raw tarball bytes---not merely its listing in the
registry's version metadata---are permanently deleted from npm's CDN (Section~\ref{sec:new-incidents}). This
is narrower than a general claim that ``purged packages cannot be diffed'': the version-resolution fixes in
Section~V-C mean a resolvable predecessor \emph{can} now be located even past a purge, and the diff runs
correctly whenever both endpoints' tarballs still exist. The limitation is specific to versions whose content
has been deleted at the CDN level, which is npm's standard, permanent takedown behavior for confirmed-
malicious releases; detection for such versions instead relies on the coarser
\texttt{requested\_version\_unresolved} signal (Section~V-D), not a working capability diff of the actual
malicious payload.
\item \textbf{Drift-monitor baseline and thresholds.} The KL-divergence thresholds (0.15/0.30) and the
synthetic seed baseline (Section~VI-E) were chosen heuristically without validation against longitudinal
production history.
\item \textbf{Optional LLM verifier.} The Ollama-based second-pass verifier was not included in the automated
evaluation corpus runs in this pass, and its accuracy on ambiguous cases has not been independently
benchmarked against a labelled adversarial-README corpus.
\end{enumerate}

(The collected-but-unused maintainer-count signal is a minor, easily-closed gap rather than a structural
limitation; it is addressed once, as a future-work item, immediately below.)

\subsection{Future Work}
Priority future directions include: (i) a longitudinal field study across multiple development teams,
including a real-traffic-derived concept-drift baseline to replace the synthetic seed described in
Section~VI-E; (ii) expansion of the Shield rule set from newly discovered malicious packages, and expansion
of Sentinel's confusable-character normalization (Section~IX-E) beyond its current Cyrillic/Greek scope to
additional Unicode blocks observed in real homoglyph squats (e.g.\ Armenian, Cherokee); (iii) incorporation
of the maintainer-count signal into Sentinel's risk score once a validated threshold is established; (iv)
integration of the LLM verifier into the default scan path for high-ambiguity cases, with independent
benchmarking against a labelled adversarial-README corpus; (v) a complementary hard-fail CI gate mode; (vi)
exploring a learned classifier for Stage~2 of the Aggregator (Section~V-D)---combining affix-stripped distance,
keyboard-adjacency-weighted distance, character n-gram similarity, length delta, and download-ratio as
features---now that the two-stage structure gives such a signal a clean place to matter without being
diluted by Stage~1's dominant flag-driven conditions; and (vii) a Stage-2-only variant of the ablation study
in Section~VIII-E, invoking \texttt{Aggregator.\_stage2\_score} directly and bypassing
\texttt{\_stage1\_gates} entirely, to isolate what weighted pillar scoring alone contributes on the subset of
records that no Stage~1 condition would otherwise decide---the current ablation, by design, mainly measures
how often Stage~1 fires rather than each pillar's standalone discriminative quality once it does not.

\section{Conclusion}
We presented CIDAS, a real-time, pre-install npm package screening system that addresses four distinct
supply-chain threat classes---typosquatting, malicious install scripts, maintainer-account compromise, and
LLM package hallucination---within a single, locally-run, fail-open security framework. CIDAS's four-pillar
architecture (Contextify, Sentinel, Shield, Aggregator) combines semantic project-context awareness, registry
reputation analysis, corroborated name-similarity detection, and AST-based static scanning into explainable,
policy-aware ALLOW/WARN/BLOCK verdicts. Supporting subsystems including an HMAC-protected trust list, a
two-tier cache, an append-only audit log now carrying genuine per-pillar score history, and a
KL-divergence-based concept-drift monitor provide ongoing assurance and team-wide governance.

Evaluated across 548 automated tests and a labelled 179-package corpus against a live daemon and live
registry conditions, CIDAS achieves a pooled F1 of 0.93 (category-balanced macro-F1: 0.92), with sub-50\,ms
latency on the cache-hit path and
single-digit-second latency on the cache-miss path under typical registry conditions. A key finding of this
evaluation is that three specific, previously undetected gaps---a structural blind spot in raw edit-distance
typosquat detection, an avoidable double-tarball-download cost in cross-version diffing, and a pair of
version-resolution bugs that silently misdiagnosed 15 real 2025--2026 supply-chain incidents as a structural
diffing limitation when the true cause was fixable---accounted for a large share of the difference between
the system's original design targets and its measured behavior; all three were root-caused against real
evaluation data and closed with targeted, decoupled refinements (reputation-corroborated affix detection;
manifest-first gating; semver-ordering-aware predecessor lookup and an explicit
\texttt{requested\_version\_unresolved} Stage~1 gate) rather than parameter retuning, and all three are
validated against the same live corpus used to discover them. The version-resolution investigation
(Section~\ref{sec:new-incidents}) is also a caution against accepting a plausible-sounding root cause without
directly verifying it: the original ``purge blocks diffing'' explanation was architecturally reasonable but
false, and only direct inspection of the live registry (and of the tarball CDN itself) surfaced the real,
narrower limitation. This experience reinforces a broader methodological point consistent with the
multi-layer, shift-left supply-chain defense literature: a system's own evaluation infrastructure, run
against real conditions rather than assumed ones, is often the most direct route to finding the gaps worth
closing next---and to discovering that an assumed explanation for a gap can itself be wrong.

\section*{Acknowledgment}
The authors thank Dr. Vijayalakshmi M N (RVCE, Department of AIML) for guidance and the RVCE Experiential
Learning programme for the opportunity to develop CIDAS as an SDG~9-aligned infrastructure project.

\appendices

\section{REST API Endpoint Reference}
Table~\ref{tab:api} lists all nine CIDAS REST API endpoints.

\begin{table}[h]
\caption{CIDAS REST API Endpoints}
\label{tab:api}
\centering
\begin{tabular}{@{}p{0.42\linewidth}p{0.38\linewidth}c@{}}
\toprule
\textbf{Endpoint} & \textbf{Description} & \textbf{Auth} \\
\midrule
GET /health              & Daemon status \& drift level      & None \\
POST /scan                & Analyze a package                 & Bearer \\
POST /trust                & Add a trust-list entry            & Bearer \\
GET /trust/verify         & Audit all trust-list HMACs        & Bearer \\
DELETE /cache              & Purge expired cache entries       & Bearer \\
POST /cache/invalidate     & Force re-scan on next request     & Bearer \\
GET /audit                & Query the audit log               & Bearer \\
POST /audit/override       & Record a manual override          & Bearer \\
GET /policy                & Return resolved policy file       & Bearer \\
\bottomrule
\end{tabular}
\end{table}

\section{Policy Schema Reference}
\begin{itemize}
\item \texttt{block\_list} (array of strings): Package names always BLOCKED.
\item \texttt{trust\_list} (array of strings): Package names that bypass all pillar analysis.
\item \texttt{min\_monthly\_downloads} (integer, default 0): Reject packages below this download threshold.
\item \texttt{require\_repository\_link} (boolean, default false): Require a repository field in
\texttt{package.json}.
\item \texttt{max\_sentinel\_distance} (integer, optional): Override the default Levenshtein typosquat
distance threshold (default 2) for this project.
\item \texttt{contextify\_weight} (float, 0.0--0.5, optional): Per-project override of the Contextify pillar
weight; the remainder is split between Sentinel and Shield in their existing ratio.
\item \texttt{warn\_requires\_confirmation} (boolean, default false): Force interactive confirmation for WARN
verdicts; non-interactive contexts are treated as BLOCK.
\end{itemize}

\section{Admin Configuration Reference}
The following per-machine flags (\texttt{\textasciitilde/.cidas/config.json}) govern optional or
recently-added behavior; all default to \texttt{true} unless noted, and can be disabled without a code change:
\begin{itemize}
\item \texttt{bypass\_disabled} (boolean): When true, the npm shim refuses a bypass environment variable and
exits with a nonzero code, for CI enforcement.
\item \texttt{package\_file\_scan} (boolean, default true): When false, Shield skips downloading and
extracting the package tarball for the primary file scan.
\item \texttt{contextify\_weight} (float, 0.0--0.5): Per-machine override of the Contextify pillar weight.
\item \texttt{shield\_manifest\_gating} (boolean, default true): Governs the manifest-first gate described in
Section~V-C.
\item \texttt{typosquat\_affix\_canonicalization} (boolean, default true): Governs affix-based typosquat
detection described in Section~V-B.
\item \texttt{typosquat\_reputation\_corroboration} (boolean, default true): Governs reputation-disparity
corroboration described in Section~V-B.
\end{itemize}

\begin{thebibliography}{21}

\bibitem{duan2021} R. Duan, O. Alrawi, R. P. Kasturi, R. Elder, B. Saltaformaggio, and W. Lee, ``Towards
measuring supply chain attacks on package managers for interpreted languages,'' in \emph{Proc. 28th Annu.
Netw. Distrib. Syst. Secur. Symp. (NDSS)}, San Diego, CA, USA, Feb. 2021.

\bibitem{zimmermann2019} M. Zimmermann, C.-A. Staicu, C. Tenny, and M. Pradel, ``Small world with high
risks: A study of security threats in the npm ecosystem,'' in \emph{Proc. 28th USENIX Secur. Symp.}, Santa
Clara, CA, USA, Aug. 2019, pp. 995--1010.

\bibitem{ohm2020} M. Ohm, H. Plate, A. Sykosch, and M. Meier, ``Backstabber's knife collection: A review of
open source software supply chain attacks,'' in \emph{Proc. Detection Intrusions Malware, Vulnerability
Assessment (DIMVA)}, ser. LNCS, vol. 12223, 2020, pp. 23--43.

\bibitem{vu2020} D.-L. Vu, I. Pashchenko, F. Massacci, H. Plate, and A. Sabetta, ``Typosquatting and
combosquatting attacks on the Python ecosystem,'' in \emph{Proc. IEEE Eur. Symp. Secur. Privacy Workshops
(EuroS\&PW)}, Genoa, Italy, Sep. 2020, pp. 509--514.

\bibitem{zapata2018} R. E. Zapata, R. G. Kula, B. Chinthanet, T. Ishio, K. Matsumoto, and A. Ihara,
``Towards smoother library migrations: A look at vulnerable dependency migrations at function level for npm
JavaScript packages,'' in \emph{Proc. IEEE Int. Conf. Softw. Maintenance Evol. (ICSME)}, Madrid, Spain, Sep.
2018, pp. 559--563.

\bibitem{zahan2022} N. Zahan, T. Zimmermann, P. Godefroid, B. Murphy, C. Maddila, and L. Williams, ``What
are weak links in the npm supply chain?'' in \emph{Proc. IEEE/ACM 44th Int. Conf. Softw. Eng.: Softw. Eng.
Practice (ICSE-SEIP)}, Pittsburgh, PA, USA, May 2022, pp. 331--340.

\bibitem{vu2021} D.-L. Vu et al., ``On the feasibility of detecting software supply chain attacks,'' in
\emph{Proc. IEEE Mil. Commun. Conf. (MILCOM)}, San Diego, CA, USA, Nov. 2021.

\bibitem{reimers2019} N. Reimers and I. Gurevych, ``Sentence-BERT: Sentence embeddings using siamese
BERT-networks,'' in \emph{Proc. 2019 Conf. Empirical Methods Nat. Lang. Process. (EMNLP-IJCNLP)}, Hong Kong,
China, Nov. 2019, pp. 3982--3992.

\bibitem{spracklen2025} J. Spracklen, R. Wijewickrama, A. H. M. N. Sakib, A. Maiti, B. Viswanath, and M.
Jadliwala, ``We have a package for you! A comprehensive analysis of package hallucinations by code
generating LLMs,'' in \emph{Proc. 34th USENIX Secur. Symp. (USENIX Security 25)}, Seattle, WA, USA, Aug.
2025, pp. 3687--3706.

\bibitem{greshake2023} K. Greshake, S. Abdelnabi, S. Mishra, C. Endres, T. Holz, and M. Fritz, ``Not what
you've signed up for: Compromising real-world LLM-integrated applications with indirect prompt injection,''
in \emph{Proc. 16th ACM Workshop Artif. Intell. Secur. (AISec '23)}, Copenhagen, Denmark, Nov. 2023, pp.
79--90.

\bibitem{latendresse2022} J. Latendresse et al., ``An empirical study of account takeovers in the npm
supply chain,'' in \emph{Proc. IEEE/ACM Int. Conf. Mining Softw. Repositories (MSR)}, Pittsburgh, PA, USA,
May 2022.

\bibitem{ladisa2023} P. Ladisa, H. Plate, M. Martinez, and O. Barais, ``SoK: Taxonomy of attacks on
open-source software supply chains,'' in \emph{Proc. 2023 IEEE Symp. Secur. Privacy (S\&P)}, San Francisco,
CA, USA, May 2023, pp. 1509--1526.

\bibitem{vu2023} D.-L. Vu, Z. Newman, and J. S. Meyers, ``Bad snakes: Understanding and improving Python
package index malware scanning,'' in \emph{Proc. 45th IEEE/ACM Int. Conf. Softw. Eng. (ICSE)}, Melbourne,
Australia, May 2023, pp. 499--511.

\bibitem{feng2020} Z. Feng et al., ``CodeBERT: A pre-trained model for programming and natural languages,''
in \emph{Findings ACL: EMNLP 2020}, Online, Nov. 2020, pp. 1536--1547.

\bibitem{cappos2008} J. Cappos, J. Samuel, S. Baker, and J. H. Hartman, ``A look in the mirror: Attacks on
package managers,'' in \emph{Proc. 15th ACM Conf. Comput. Commun. Secur. (CCS '08)}, Alexandria, VA, USA,
Oct. 2008, pp. 565--574.

\bibitem{sejfia2022} A. Sejfia and M. Sch\"{a}fer, ``Practical automated detection of malicious npm
packages,'' in \emph{Proc. 44th IEEE/ACM Int. Conf. Softw. Eng. (ICSE '22)}, Pittsburgh, PA, USA, May 2022,
pp. 1681--1692.

\bibitem{ferreira2021} G. Ferreira, L. Jia, J. Sunshine, and C. K\"{a}stner, ``Containing malicious package
updates in npm with a lightweight permission system,'' in \emph{Proc. 43rd IEEE/ACM Int. Conf. Softw. Eng.
(ICSE)}, Madrid, Spain, May 2021, pp. 1334--1346.

\bibitem{froh2023} F. N. Froh, M. F. Gobbi, and J. Kinder, ``Differential static analysis for detecting
malicious updates to open source packages,'' in \emph{Proc. ACM Workshop Softw. Supply Chain Offensive Res.
Ecosystem Defenses (SCORED '23)}, Copenhagen, Denmark, Nov. 2023, pp. 41--49.

\bibitem{birsan2021} A. Birsan, ``Dependency Confusion: How I Hacked Into Apple, Microsoft and Dozens of
Other Companies,'' \emph{Medium}, Feb. 9, 2021.

\end{thebibliography}

\end{document}

